\scanpage{187}
des $G_{\mathbb{R}}$-modules de dim finie sur $\mathbb{R}$. Le cône qui exprime cette effectivité est \uline{engendré} par les \deleted{$\unsure{u(s')|v(s')}$} rayons définis par les $v(s')$. \unsure{C'est} un fait qui exprime \inserted{sous forme ``analytique''} le fait que les \unsure{frobénius|frobenius|Frobenius} $f(s')\in G_{\mathbb{R}}^{\xi(s')}(\mathbb{R})$ ont la vertu que leurs valeurs propres, dans \unsure{tte} représentation de $G_{\mathbb{R}}^{\xi(s')}\deleted{(\mathbb{R})}$ -- et en fait \unsure{$=$|même} de $G_{\deleted{\mathbb{R}}\mathbb{C}}^{\xi(s')}$ -- sont des \uline{entiers} algébriques, ce qui est une condition arithmétique assez draconienne !
% archive sheet number: 186
